% -----------------------------------------------------------------------------
% Beginner bridge: how ordinary one-particle quantum mechanics emerges
% -----------------------------------------------------------------------------

量子场算符可以改变粒子数，而普通量子力学的波函数只是复数值函数；两种语言之间的联系最好先在一个守粒子数的非相对论 Schr\"odinger 场中看清。这个低能模型把 Fock 空间按粒子数分解，并使普通单粒子量子力学直接作为其中的固定一粒子扇区出现。

等时正则关系为
\begin{equation}
 [\hat\Psi(\mathbf r),\hat\Psi^\dagger(\mathbf r')]
 =\delta^{(3)}(\mathbf r-\mathbf r'),
 \qquad
 [\hat\Psi,\hat\Psi]=0.
 \label{eq:sch-field-comm-r6}
\end{equation}
若真空满足 $\hat\Psi(\mathbf r)|0\rangle=0$，则 $\hat\Psi^\dagger(\mathbf r)|0\rangle$ 是在位置标签 $\mathbf r$ 上产生一个粒子的广义态。任意一粒子态因此可展开为
\begin{equation}
 \boxed{
 |\Psi_1(t)\rangle
 =\int\dd^3 r\,
 \psi_{\rm 1p}(\mathbf r,t)
 \hat\Psi^\dagger(\mathbf r)|0\rangle .}
 \label{eq:fock-oneparticle-expansion-dp68}
\end{equation}
这里的 $\psi_{\rm 1p}$ 才是普通第一量子化中的复数值波函数；$\hat\Psi$ 本身仍是作用在 Fock 空间上的算符值分布。

对自由哈密顿量
\begin{equation*}
 \hat H
 =\int\dd^3 r\,
 \hat\Psi^\dagger(\mathbf r)
 \left(-\frac{\hbar^2\nabla^2}{2m}\right)
 \hat\Psi(\mathbf r),
\end{equation*}
Schr\"odinger 绘景的态方程在一粒子扇区中约化为
\begin{equation}
 \boxed{
 \ii\hbar\partial_t\psi_{\rm 1p}(\mathbf r,t)
 =-\frac{\hbar^2}{2m}\nabla^2\psi_{\rm 1p}(\mathbf r,t).}
 \label{eq:fock-oneparticle-schrodinger-dp68}
\end{equation}
\eqnrefs{eq:fock-oneparticle-schrodinger-dp68} 不是另加的一条动力学假设，而是把场论哈密顿量限制到\eqnrefs{eq:fock-oneparticle-expansion-dp68} 所定义的一粒子扇区后的结果。

$\hat\Psi^\dagger$ 用来构造粒子态，$\psi_{\rm 1p}$ 用来描述固定一粒子态在位置基中的展开系数。进入两粒子及更高扇区后，态由多个产生算符构造，玻色对易关系或费米反对易关系自动决定交换对称性，因此不需要再把场算符误解成一条“多粒子波函数”。


固定粒子数扇区解释了普通波函数从哪里来。下一步回到相对论场论：此时粒子数不再必须守恒，真正需要保留的是局域场及其二点函数；类空分离处的场对易关系将给出微观因果性。

